\chapter{Downstream Inference}
\label{ch:downstream}

Chapters~\ref{ch:serial} and~\ref{ch:cross} deliver the two objects the rank-inference
procedure consumes: the pairwise HAC standard errors $\seNW{jk}$ and a positive
semi-definite covariance ($\Sh$ from the MDS reconstruction or $\Omegahat$ from the
direct construction). The procedure of \citet{mogstad2024}, used unchanged, turns them
into simultaneous rank confidence sets, and a single implementation serves both
covariance constructions. The chapter closes with the companion $\tau$-best confidence
sets.

\section{Simulating the studentized-max null}
\label{sec:studentized}

Ranking reduces to simultaneous inference on the signs of the contrasts
$\Delta_{j,k}=\theta_j-\theta_k$ (\S\ref{sec:ranks}). For an ordered pair $(j,k)$ the
one-sided studentized statistic
\begin{equation}
  T_{j,k} \;=\; \frac{\hat{\Delta}_{j,k}}{\seNW{jk}}
\end{equation}
is large and positive when $\theta_j>\theta_k$ is well established, that is, when $k$ is
confirmed better than $j$. Controlling the maximum of these statistics over the family
of active pairs controls the family-wise error rate of all the sign decisions at once.

The critical value is the $(1-\alpha)$ quantile of that maximum under the null
$\Delta_{j,k}=0$. I obtain it by simulation: draw $B$ Gaussian vectors from the PSD
covariance, form the studentized contrast for each active pair, and take the maximum.
With $A$ the current set of active ordered pairs and $D^{(b)}$ the $b$-th simulated
contrast vector,
\begin{equation}
  T^{(b)} \;=\; \max_{(j,k)\in A}\frac{D^{(b)}_{j,k}}{\seNW{jk}},
  \qquad
  \hat{c}(1-\alpha, A) \;=\; \widehat{Q}_{1-\alpha}\big(\{T^{(b)}\}_{b=1}^{B}\big),
  \label{eq:cv-sim}
\end{equation}
where $D^{(b)}_{j,k}=Z^{(b)}_j-Z^{(b)}_k$ for the MDS draw $Z^{(b)}\sim\Normal(0,\Sh)$,
or $D^{(b)}$ is drawn directly from $\Normal(0,\Omegahat)$ (\S\ref{sec:seam}). A pair is
confirmed ($\theta_j>\theta_k$, so $k$ ranks above $j$) when
$\hat{\Delta}_{j,k} - \hat{c}\,\seNW{jk} > 0$.

This simulation is a \emph{parametric bootstrap from the estimated covariance}. It
approximates the unknown null distribution of the maximum by drawing from the Gaussian
law the contrast estimators follow asymptotically, with the covariance replaced by its
estimate $\Sh$ or $\Omegahat$, rather than by resampling the data. The validity of the
critical value therefore rests on that estimate, which is why Chapters~\ref{ch:serial}
and~\ref{ch:cross} are devoted to it. I call it the \emph{simulation route} below.

\paragraph{A nonparametric benchmark.}
\label{par:mbb}
Resampling the data is another way to obtain the critical value, without any covariance
estimate. I use the block bootstrap of \citet{kunsch1989} in the circular form of
\citet{politisromano1992}. Each replication draws one set of $n$ row indices for the
whole panel, as contiguous blocks of $\ell$ periods that wrap around from the last
period to the first, so every period is equally likely to be drawn. On the rows of that
resample where both $j$ and $k$ are observed, it recomputes the contrast
$\hat{\Delta}^{*}_{j,k}$; then $D^{(b)}_{j,k}=\hat{\Delta}^{*}_{j,k}-\hat{\Delta}_{j,k}$
takes the place of the Gaussian draw in \eqref{eq:cv-sim}, with $\seNW{jk}$ kept fixed.
Joint resampling of rows carries the cross-sectional dependence, blocks carry the serial
dependence, and the missingness pattern travels with the rows. The block bootstrap
variance of a mean with block length $\ell$ is asymptotically the Bartlett estimator
with bandwidth $\ell-1$, so I set $\ell=L+1$, with $L$ the median Andrews bandwidth of
\S\ref{sec:bandwidth} across pairs. Because rows are resampled jointly, one block length
serves all pairs, while each standard error uses its own pair's bandwidth. The bootstrap
contrasts are therefore close to unit variance only for pairs whose bandwidth is close to
the median. The two routes also differ in how cross-sectional dependence enters
(estimated in $\Omegahat$ or implicit in the resample) and in the shape of the null
distribution (Gaussian or resampled). The resamples are drawn once and reused in every
stepdown round, as the Gaussian draws are (\S\ref{sec:seam}). I report the bootstrap as
a check on the simulation route, not as a competing method, since it bypasses the object
this thesis estimates.

\section{The Romano--Wolf stepdown}
\label{sec:stepdown}

A single application of \eqref{eq:cv-sim} over all pairs is valid but conservative. The
stepdown of \citet{romanowolf2005} recovers power while retaining family-wise control:

\begin{enumerate}[nosep]
  \item Initialise the active set $A$ to all ordered pairs (with finite $\seNW{jk}$).
  \item Compute $\hat{c}(1-\alpha, A)$ from \eqref{eq:cv-sim} and reject every active
        pair with $\hat{\Delta}_{j,k} - \hat{c}\,\seNW{jk} > 0$.
  \item If any pair was rejected, remove the rejected pairs from $A$ and return to
        step~2; otherwise stop.
\end{enumerate}

Because the maximum in \eqref{eq:cv-sim} is taken over a shrinking set, the critical
values are non-increasing across rounds, and each round can only confirm more pairs.
This monotone tightening is weakly more powerful than the single-step test and still
controls the family-wise error rate. At termination the confirmed pairs give the
better/worse counts $N^+_j,N^-_j$ and hence the rank confidence set
$\rankci{j}=[\,N^+_j+1,\;p-N^-_j\,]$ of \eqref{eq:rankset}. \emph{Marginal}
(per-forecaster) confidence sets follow by restricting the family to the $2(p-1)$ pairs
involving a single $j$ and computing its own critical value. They are tighter, but they
control coverage for one forecaster at a time rather than jointly.

\section{The covariance seam}
\label{sec:seam}

The two constructions of Chapter~\ref{ch:cross} live in different spaces. MDS returns a
$p\times p$ level covariance $\Sh$, from which one draws $Z\in\Real^p$ and then
differences. The direct construction returns a $q\times q$ difference covariance
$\Omegahat$, from which one draws $D\in\Real^q$ directly. So that a single stepdown can
use either, the interface returns not a raw covariance but the \emph{studentized
difference draws}: a $B\times q$ array $T$ with, for unordered pair $a=(j,k)$,
\begin{equation}
  \text{(MDS)}\quad T^{(b)}_a=\frac{Z^{(b)}_j-Z^{(b)}_k}{\seNW{a}},
  \qquad
  \text{($\OmegaDirect$)}\quad T^{(b)}_a=\frac{D^{(b)}_a}{\sqrt{\Omegahat_{aa}}},
  \quad Z^{(b)}\sim\Normal(0,\Sh),\ D^{(b)}\sim\Normal(0,\Omegahat).
\end{equation}
The two routes studentize differently. MDS divides by the Newey--West standard error
$\seNW{a}$, the quantity to which $\Sh$ is calibrated. The direct construction divides
by $\sqrt{\Omegahat_{aa}}$, the diagonal of the projected matrix it draws from, so its
draws have unit variance by construction; the same quantity replaces $\seNW{a}$ in its
observed statistics and in the rejection rule of \S\ref{sec:studentized}. On a
contiguous support the unprojected diagonal equals $(\seNW{a})^2$
(\S\ref{sec:omega-nesting}). The projection can only raise it, so the direct route's
standard errors are at least $\seNW{a}$, and the two coincide when the projection leaves
the diagonal unchanged. For MDS, by contrast, the projection can change the variance of
the simulated contrasts, so its draws need not have unit variance. They are weakly
over-dispersed, since the projection only adds variance; Chapter~\ref{ch:mc} shows the
consequences.

The draws are taken once, and the stepdown restricts them to the active columns in each
round. This is both the monotone procedure and cheaper than re-drawing. A single runtime
switch, $\text{covariance}\in\{\textsc{mds},\OmegaDirect\}$, selects the construction.
Everything downstream is identical, which keeps the two approaches cleanly comparable
(Chapter~\ref{ch:mc}) and independently removable.

\section{Confidence sets for the $\tau$-best}
\label{sec:taubest}

Beyond the rank of each forecaster, one often wants the set of forecasters that could
belong to the top $\tau$: the $\tau$-\emph{best} confidence set of
\citet[Section~3.4]{mogstad2024}. A (random) set $\setbest\subseteq\{1,\dots,p\}$ is a
$\tau$-best confidence set at level $1-\alpha$ if
\begin{equation}
  \Prob\big(\text{every truly top-$\tau$ forecaster lies in } \setbest\big)\;\ge\;1-\alpha.
  \label{eq:tau-coverage}
\end{equation}

\paragraph{Naive projection.} Given the joint rank confidence sets
$\rankci{j}=[L_j,U_j]$ from \S\ref{sec:stepdown}, the projection
\begin{equation}
  {\setbest}^{\text{naive}}=\{\, j : L_j\le\tau\,\}
\end{equation}
inherits coverage \eqref{eq:tau-coverage} by elementary set inclusion. With probability
at least $1-\alpha$ every $\rankci{j}$ covers its forecaster's true rank $r_j$. On that
event a truly top-$\tau$ forecaster has $r_j\le\tau$, hence $L_j\le r_j\le\tau$, and so
lies in ${\setbest}^{\text{naive}}$. Only the lower bound matters: a forecaster whose
whole set lies within the top $\tau$ ($U_j\le\tau$) is top-$\tau$ for certain on that
event, and belongs in the set just as much as one whose set straddles $\tau$. The
projection is conservative in general and comes free with the rank confidence sets.

\paragraph{Direct procedure.} A dedicated test statistic gives a sharper set. For
forecaster $j$,
\begin{equation}
  T_{n,j}\;=\;\min_{K\in\mathcal{K}}\;\max_{k\notin K\cup\{j\}}\;
  \frac{\hat{\Delta}_{j,k}}{\seNW{jk}},
  \qquad
  \mathcal{K}=\big\{K\subseteq\{1,\dots,p\}:\;|K|=\tau-1\big\},
  \label{eq:tau-stat}
\end{equation}
where $K$ exempts $\tau-1$ forecasters (candidates for the other top slots), the inner
$\max$ is the strongest evidence that some non-exempt $k$ beats $j$, and the outer
$\min$ grants $j$ the most favourable exemption. A large $T_{n,j}$ means that no matter
which $\tau-1$ forecasters are granted the other top slots, someone outside is still
substantially better than $j$, so $j$ is unlikely to be top-$\tau$.

\begin{remark}[Sign convention]
\citet{mogstad2024} write the procedure for ``larger $\theta$ is better'' (e.g.\
achievement scores). For MSE-style losses ``smaller $\theta$ is better'', so the
procedure is applied to $-\theta$. By their Remark~3.11 this flips the numerator:
$(\hat\theta_k-\hat\theta_j)/\seNW{jk}$ becomes
$\hat{\Delta}_{j,k}/\seNW{jk}$ as in \eqref{eq:tau-stat}. Large positive values still
mean ``$j$ has a much larger loss than some non-exempt $k$''.
\end{remark}

\paragraph{Critical value and stepwise testing.} The null $H_j$ (``$j$ is among the top
$\tau$'') is rejected, and $j$ excluded from $\setbest$, when
$T_{n,j}>\hat{c}_n(1-\alpha,I_n)$. The critical value is the largest per-exemption
quantile,
\begin{equation}
  \hat{c}_n(1-\alpha,I_n)=\max_{K\in\mathcal{K}}
  \widehat{Q}_{1-\alpha}\!\Big(\big\{\,\textstyle\max_{j\in I_n}\max_{k\notin K\cup\{j\}}
  D^{(b)}_{j,k}/\seNW{jk}\,\big\}_{b}\Big),
\end{equation}
where the null contrasts $D^{(b)}$ are drawn from the MDS covariance as in
\S\ref{sec:studentized}. The set is built by the same stepdown logic as
\S\ref{sec:stepdown}: initialise $I_n=\{1,\dots,p\}$, exclude every $j$ with
$T_{n,j}>\hat{c}_n$, shrink $I_n$, and repeat until no further exclusions; the survivors
are $\setbest$. As before, removing excluded forecasters cannot raise the critical value,
so the stepwise refinement is weakly more powerful than a single step. The reported
results use the projection, not the direct procedure.
